\section{An optional weakening: transport of two margin levels}
\label{sec:two-margin-transport}

The statistical argument needs separation of the training and test
margin levels, rather than stability of every intermediate score.
The following optional prior condition expresses precisely that
separation. It is a mathematical extension, not the condition used in
the primary empirical protocol.

For each network use the same calibrated radius
$r_c(w)=c/\sqrt{1+H E_v(w)}$ as before. Require, for every input $x$,
every perturbation with $\|\eta\|_*\le r_c(w)$, and both
$y\in\{-1,1\}$, the implications
\begin{align}
 yz_w(x)\ge\tfrac12
 &\quad\Longrightarrow\quad yz_w(x;\eta)\ge\tfrac7{16},
 \label{eq:transport-training-level}\\
 yz_w(x)\le\tfrac14
 &\quad\Longrightarrow\quad yz_w(x;\eta)\le\tfrac5{16}.
 \label{eq:transport-test-level}
\end{align}
Because both choices of $y$ are required, this is a label-free condition
on the network. It asks a perturbation to preserve a gap between a
confident training output and an output counted as a test error.
The condition is Borel measurable: extrema over a compact perturbation
ball are continuous functions of the weights, and the Boolean input
set and label set are finite.

If the original global clean-margin condition $|z_w(x)|>1/2$ is
also retained, the second implication is redundant. The first one
for the network's predicted sign guarantees perturbed margin at least
$7/16$; negating this inequality gives the second implication for the
opposite label, the only label for which it is triggered. In this original
globally separated model, the restriction is an ordinary all-layer
adversarial margin requirement. The second implication is needed when
the theorem permits globally uncertain outputs.

\begin{proposition}[The same bound under two-margin transport]
\label{prop:two-margin-transport}
Replace the clipped-score support restriction in the main theorem by
\eqref{eq:transport-training-level}--\eqref{eq:transport-test-level}.
All its sample bounds and all-width conclusions remain valid, with the
same training likelihood, test margin, and posterior-confidence
consequences. The same replacement is valid for the localized versions,
where the implications are imposed outside the permitted exception set.
\end{proposition}

\begin{proof}
Use the same common compression radius $r\le1/6$. On a successful
compression coupling, the candidate's raw output differs from a
perturbed original output by at most $r/8\le1/48$.
At an observed training input, $yz_w(x)>1/2$, so
\eqref{eq:transport-training-level} implies
\[
 yh(x)\ge\frac7{16}-\frac1{48}=\frac5{12}.
\]
At a test input where $yz_w(x)\le1/4$,
\eqref{eq:transport-test-level} implies
\[
 yh(x)\le\frac5{16}+\frac1{48}=\frac13.
\]
Thus the same ramp, zero at and above $5/12$ and one at and below
$1/3$, has zero training loss on successful couplings and dominates
the original test-error indicator on successful couplings. The
compression entropy and readout-tail bounds are unchanged.

Every weight in the target cylinder has the correct signed output
greater than one both before and after every admissible perturbation.
It therefore satisfies both implications for both labels. The original
raw Gaussian target-mass lower bound, and hence the posterior entropy
bound, remain valid. The rest of the proof is identical. Masking by the
controlled exception sets gives the localized version.
\end{proof}

The margin-band score
$U(z)=\operatorname{sign}(z)\min\{1/4,(|z|-1/4)_+\}$ gives a
sufficient condition for these implications: if
$|U(z_w(x;\eta))-U(z_w(x))|\le1/16$, then a clean signed score at
least $1/2$ has perturbed $U$-score at least $3/16$, hence perturbed
signed score at least $7/16$. A clean signed score at most $1/4$ has
$U$-score at most zero, so its perturbed signed score is at most
$5/16$. Oddness of $U$ handles either label.

The implications allow additional changes at intermediate scores.
For example, the pair of scores $3/8$ and $11/20$ respects both
implications for both labels, although its $U$-score change is $1/8$.
Such an increase in confidence need not affect either of the two
loss-transfer steps. This explains the weakening; it does not establish
that trained networks satisfy it globally.

More generally, with training margin $m$, test margin $u<m$, and a
noise tolerance $a<(m-u)/2$, replace the four numerical levels by
$m,m-a,u,u+a$. If the compressor's output-rounding error is $e$ and
$m-u-2a-2e>0$, the same proof uses a ramp across this positive gap.
Its sample dependence is polynomial in the inverse gap. As before,
an empirical assessment of this optional condition would require a
separate stated protocol; the primary experiments are not relabeled.
